\clearpage\fancyhead[L]{\small PALLAS.}

\begin{center}\textbf{PALLAS.}\par\vspace{4pt}\hrule\vspace{4pt}\fontsize{10}{12}\selectfont Göttingische gelehrte Anzeigen. Stück 32. S. 305\,..\,308. 24. Februar 1810.\par\vspace{4pt}\hrule\vspace{4pt}\end{center}

\par\indent \emph{Pallas} war im Jahr 1808 nur an wenigen Sternwarten und nur unvollkommen beobachtet worden. Deshalb hatte man die Elemente seit 1807 auch nicht berichtigen können. An der hiesigen Sternwarte machte die große Lichtschwäche des Planeten jede Beobachtung am Mauerquadranten unmöglich. Die wenigen Beobachtungen mit dem Kreismikrometer konnten wegen besonderer Umstände keine große Genauigkeit beanspruchen. An der Seeberger Sternwarte beobachtete Herr \textsc{von Lindenau} den Planeten zwar mit gewohnter Genauigkeit am Passageninstrument, doch konnten auch dort keine Deklinationen beobachtet werden. Die Bestimmungen mit dem Äquatorialsektor in Mailand setzen erst am 22. August ein und waren daher für die Opposition kaum noch zu gebrauchen. Erst durch das Astronomische Jahrbuch für 1812 sind nun noch einige vollständige Beobachtungen bekannt geworden, die Herr \textsc{David} um die Zeit der Opposition in Prag gemacht hat. Professor \textsc{Gauss} hat deren Deklinationen zusammen mit den an der Seeberger Sternwarte bestimmten Rektaszensionen verwendet, um die Opposition abzuleiten. Leider hat die Untersuchung der Prager Beobachtungen gezeigt, dass sie schlecht zusammenpassen. So ergeben die verschiedenen Deklinationen zum Beispiel Resultate, die fast eine Minute voneinander abweichen. Da jedoch keine besseren Deklinationen vorliegen, blieb nur, die vorhandenen so gut wie möglich zu nutzen. So ergab sich für die Opposition von 1808, die fünfte bisher beobachtete, folgendes Ergebnis.

\begin{center}Opposition der Pallas am 26. Juli 1808.\\$21^{\text{u}}\;17^{\text{m}}\;32^{\text{s}}$ Meridian von Göttingen\par Wahre Länge .\;.\;.\;.\;.\; $304^\circ\;2'\;59''7$\par Wahre geozentrische Breite\quad $37\;33\;54$ nördlich\end{center}

\Needspace{26\baselineskip}\SourceLeaf{323}\par\indent Im Jahr 1809 scheinen die Beobachtungen der Pallas an der hiesigen Sternwarte (siehe Gel. Anz. 1809, S. 1348) [26. Aug. 1809, S. 314 dieses Bandes] die frühesten gewesen zu sein. Von anderen Sternwarten hat Professor \textsc{Gauss} durch die freundliche Mitteilung der Herren \textsc{von Lindenau} und \textsc{Bouvard} eine wertvolle Reihe von Seeberger und Pariser Meridianbeobachtungen erhalten. Die Pariser Beobachtungen, die sonst erst nach einigen Jahren bekannt würden, geben wir im Folgenden wieder. Die mittleren Zeiten, die Herr \textsc{Bouvard} nicht mitgeteilt hatte, hat Professor \textsc{Gauss} selbst ergänzt.

\Needspace{18\baselineskip}\begin{center}\emph{Beobachtungen der Pallas in Paris.}\end{center}
\begin{center}\fontsize{9}{11}\selectfont\setlength{\tabcolsep}{5pt}\renewcommand{\arraystretch}{1.0}\begin{tabular}{lrrr}
1809 & Mittlere Zeit & \shortstack{Scheinbare\\Rektaszension} & \shortstack{Scheinbare südliche\\Deklination}\\
Sept. 12 & $12^{\text{u}}\;51^{\text{m}}\;39^{\text{s}}6$ & $4^\circ\;30'\;25''80$ & $4^\circ\;20'\;10''4$\\
15 & $12\;37\;41.7$ & $3\;57\;48.75$ & $5\;5\;16.4$\\
28 & $11\;36\;52.6$ & $1\;31\;50.25$ & $8\;20\;23.4$\\
29 & $11\;32\;11.1$ & $1\;20\;24.30$ & $8\;34\;57.3$\\
Okt. 2 & $11\;18\;7.1$ & $0\;46\;15.30$ & $9\;18\;22.5$\\
5 & $11\;4\;5.9$ & $0\;12\;35.25$ & $10\;0\;21.0$\\
6 & $10\;59\;25.5$ & $0\;1\;37.50$ & $10\;14\;6.9$\\
7 & $10\;54\;46.2$ & $359\;50\;44.55$ & $10\;27\;36.8$\\
8 & $10\;50\;7.1$ & $359\;39\;54.90$ & $10\;40\;1.5$\\
9 & $10\;45\;28.7$ & $359\;29\;16.50$ & $10\;54\;5.5$\\
13 & $10\;27\;2.2$ & $358\;48\;26.25$ & $11\;54\;59.5$\\
\end{tabular}\end{center}

\par\indent Die hier abgedruckten Zahlen stimmen genau mit den Angaben im Brief von Herrn \textsc{Bouvard} überein. Die Auswertung dieser Bestimmung hat jedoch gezeigt, dass die Rektaszension vom 12. September um etwa $1'\,15''$ im Bogen zu groß ist (die Rektaszensionen waren nur in Zeit angegeben), die Deklination vom 8. Oktober um 1 Minute zu klein und die Deklination vom 13. Oktober um 10 Minuten zu groß. Aus diesen Beobachtungen ergeben sich die folgenden Bestimmungen für die sechste Opposition der Pallas.

\begin{center}Zeit der Opposition am 22. Sept. 1809.\\$16^{\text{u}}\;10^{\text{m}}\;20^{\text{s}}$ Meridian von Göttingen\par Wahre Länge der Pallas .\;.\; $359^\circ\;40'\;4''4$\par Wahre geozentrische Breite ,\quad $7\;22\;10.1$ südlich\end{center}

\par\indent Die Ergebnisse, die Professor \textsc{Gauss} aus diesen beiden Oppositionen in Verbindung mit den vier früheren abgeleitet hat, werden bei anderer Gelegenheit bekannt gemacht.

Wie im vorigen Jahr drucken wir hier noch ein Verzeichnis ab, das uns kürzlich zur Veröffentlichung zugesandt wurde. Es enthält einige im Jahr 1810 sichtbare Sternbedeckungen, die in den astronomischen Ephemeriden entweder gar nicht oder nicht richtig angegeben sind. Herr \textsc{Wisniewsky} hat sie für den Meridian und Horizont von Berlin vorausberechnet. --- ---
